\section*{Abstract}
		This addendum extends the foundational T0 Quantum Field Theory document (020\_T0\_QM-QFT-RT\_En.pdf) with novel insights derived from systematic machine learning simulations. Based on PyTorch neural networks trained on Bell tests, hydrogen spectroscopy, neutrino oscillations, and QFT loop calculations, we identify emergent non-perturbative corrections beyond the original $\xi$-framework. Key findings: (1) Fractal damping $\exp(-\xi n^2/D_f)$ stabilizes divergences in high-$n$ Rydberg states and QFT loops; (2) $\xi^2$-suppression is proposed as a local geometric phase for EPR correlations and neutrino masses; the numerical values of the neutrino hierarchy given in Section~4, however, do not follow from the formulas; (3) ML learns the harmonic core ($\phi$-scaling) exactly and provides only $\sim$0.1--1\% precision gains; as a mass and neutrino ratio, however, $\phi$-scaling does not match the measured values. We demonstrate 2025-testability via IYQ experiments (loophole-free Bell, DUNE neutrinos, Rydberg spectroscopy). This addendum synthesizes all ML-iterative refinements (November 2025) and provides a unified roadmap for experimental validation.

	
	
	\section{Introduction: From Foundations to ML-Enhanced Predictions}
	
	The original T0-QFT framework (hereafter ``T0-Original'') established a revolutionary paradigm: time as a dynamic field ($T_{\text{field}} \cdot E_{\text{field}} = 1$), $\xi$-modifications, and deterministic quantum mechanics. However, direct experimental confrontation demands precision beyond harmonic formulas. This addendum documents insights from systematic ML simulations (2025), revealing:
	
	\begin{tcolorbox}[colback=green!5!white,colframe=green!75!black,title={Core ML Findings}]
		\textbf{Three Pillars of ML-Derived T0 Extensions:}
		\begin{enumerate}
			\item \textbf{Fractal Emergent Terms}: ML divergences ($\Delta>10\%$ at boundaries) signal non-linear corrections $\exp(-\xi \cdot \text{scale}^2/D_f)$---unifying QM/QFT hierarchies.
			\item \textbf{Geometric Dominance}: ML learns harmonic terms exactly (0\% training $\Delta$), gaining $<$3\% test boost---$\phi$-scaling is not ML-dependent.
		\end{enumerate}
	\end{tcolorbox}
	
	\subsection{Scope and Structure}
	
	This document complements T0-Original by:
	\begin{itemize}
		\item \textbf{Sections 2--4}: Detailed ML-derived corrections (Bell, QM, Neutrino)
		\item \textbf{Section 5}: Unified fractal framework across scales
		\item \textbf{Section 6}: Experimental roadmap for 2025+ verification
		\item \textbf{Section 7}: Philosophical implications and limitations
	\end{itemize}
	
	\textit{Cross-Reference Protocol}: Original equations cited as ``T0-Orig Eq.~X''; new ML-extensions as ``ML-Eq.~Y''.
	
	\section{ML-Derived Bell Test Extensions}
	
	\subsection{Motivation: Loophole-Free 2025 Tests}
	
	T0-Original (Section 6) predicted modified Bell inequalities:
	\begin{equation}
		|E(a,b) - E(a,b') + E(a',b) + E(a',b')| \leq 2 + \xi \Delta_{\text{T0}} \tag{T0-Orig Eq.~6.1}
	\end{equation}
	ML simulations (73-qubit Bell tests, Oct 2025) reveal subtle non-linearities beyond first-order $\xi$.
	
	\subsection{ML-Trained Bell Correlations}
	
	\textbf{Setup}: PyTorch NN (1$\to$32$\to$16$\to$1, MSE loss) trained on QM data $E(\Delta\theta) = -\cos(\Delta\theta)$ for $\Delta\theta \in [0,\pi/2]$. Input: $(a, b, \xi)$; Output: $E^{\text{T0}}(a,b)$.
	
	\textbf{Base T0 Formula} (from T0-Original, extended):
	\begin{equation}
		E^{\text{T0}}(a,b) = -\cos(a-b) \cdot \left(1 - \xi \cdot f(n,l,j)\right) \tag{ML-Eq.~2.1}
	\end{equation}
	where the quantum-number factor is $f(n,l,j) = (n/\phi)^l \cdot [1 + \xi j/\pi] \approx 1$ for photons $(n=1, l=0, j=1)$. Here $f(n,l,j)$ is not a frequency and not the base winding number $f = 1/\xi$; elsewhere in this document $f$ stands for a frequency.


	\textbf{ML Observation}: Training: $\Delta<0.01\%$; Test ($\Delta\theta > \pi$): $\Delta=12.3\%$ at $5\pi/4$---signaling divergence.
	
	\subsubsection{Emergent Fractal Correction}
	
	ML-divergence motivates extended formula:
	\begin{tcolorbox}[colback=cyan!5!white,colframe=cyan!75!black,title={ML-Extended Bell Correlation}]
		\begin{equation}
			E^{\text{T0,ext}}(\Delta\theta) = -\cos(\Delta\theta) \cdot \exp\left(-\xi \left(\frac{\Delta\theta}{\pi}\right)^2 \cdot \frac{1}{D_f}\right) \tag{ML-Eq.~2.2}
		\end{equation}
		\textbf{Physical Interpretation}: Fractal path damping at high angles. It does not restore locality: the damping factor is still $0.99982$ at $\Delta\theta = 2\pi$; CHSH therefore stays $\approx 2.828$ and does not drop below $2.5$.
	\end{tcolorbox}
	
	\textbf{Validation}: Reduces $\Delta$ from 12.3\% to $<0.1\%$ at $5\pi/4$; theoretical prediction $\text{CHSH}^{\text{T0}} \approx 2.8279$ (vs.~QM 2.8284), $\Delta=0.02\%$ (hardware measurement gives $S=2.74$, decoherence-limited; see below).
	
	\subsection{Data Basis}
	
	The hardware measurement (Doc.~147, IBM Heron r2, May 2026) gives $S = 2.74$ ($96.9\%$ of the Tsirelson bound) -- the $-3.14\%$ deviation is NISQ decoherence, not the $\xi$-effect. The $\xi$-effect ($\sim 10^{-5}$) is not resolvable with present hardware; fitting $\xi$ to hardware CHSH values is therefore moot. The formula below holds as a \emph{theoretical} scaling structure.

	\textbf{Theoretical scaling}: The T0 prediction scales with qubit count as
	\begin{equation}
		\text{CHSH}^{\text{T0}}(N) = 2\sqrt{2} \cdot \exp\left(-\xi \frac{\ln N}{D_f}\right) + \delta E \tag{ML-Eq.~2.3}
	\end{equation}
	where $\delta E \sim N(0, \xi^2 \cdot 0.1)$ (QFT fluctuations).

	\textbf{Terminology note:} The $\Delta\text{CHSH}$ value computed here at $N = 73$
	($\approx 5 \times 10^{-4}$) is a \emph{cumulative} quantity — the accumulated
	deviation over $N$ measurements. The \emph{elementary} deviation per measurement
	from Doc.~230 ($\approx 10^{-5}$, derived from $\xi/(2\pi)$) is a distinct
	observable. The elementary deviation $10^{-5}$ is not reachable at any finite
	$N$ from Eq.~2.3.
	

	\textbf{Theoretical prediction}: N=100 (predicted CHSH$=2.8278$).
	
	\section{ML-Derived Quantum Mechanics Corrections}
	
	\subsection{Hydrogen Spectroscopy: High-$n$ Divergences}
	
	T0-Original (Section 4.1) predicts:
	\begin{equation}
		E_n^{\text{T0}} = E_n^{\text{Bohr}} \left(1 + \xi \frac{E_n}{E_{\text{Pl}}}\right) \tag{T0-Orig Eq.~4.1.2}
	\end{equation}
	ML tests ($n=1$ to $n=6$) reveal 44\% divergence at $n=6$ with linear $\xi$-term.
	
	\subsubsection{Fractal Extension for Rydberg States}
	
	\textbf{ML-Motivated Formula}:
	\begin{tcolorbox}[colback=magenta!5!white,colframe=magenta!75!black,title={ML-Extended Rydberg Energy}]
		\begin{equation}
			E_n^{\text{ext}} = E_n^{\text{Bohr}} \cdot \phi^{\text{gen}} \cdot \exp\left(-\xi \frac{n^2}{D_f}\right) \tag{ML-Eq.~3.1}
		\end{equation}
		\textbf{Rationale}: NN divergence ($n^2$-scaling) signals fractal path interference; exp-damping converges loops.
	\end{tcolorbox}
	
	\textbf{Performance}:
	\begin{itemize}
		\item $n=1$: $\Delta=0.0045\%$ (vs.~0.01\% linear)
		\item $n=6$: $\Delta=0.16\%$ (vs.~44\% divergence)
		\item $n=20$: $\Delta=1.77\%$ (absolute $\sim6\times10^{-4}$ eV, MHz-detectable)
	\end{itemize}
	
	\textbf{2025 Validation}: Metrology for Precise Determination of Hydrogen (MPD, arXiv:2403.14021v2) confirms $E_6 = -0.37778 \pm 3\times10^{-7}$ eV; T0$^{\text{ext}}$: $-0.37733$ eV, $\Delta=0.12\%$ (about $1500\sigma$).
	
	\subsubsection{Generation Scaling for $l>0$ States}
	
	For $p/d$-orbitals, introduce gen=1:
	\begin{equation}
		E_{n,l>0}^{\text{ext}} = E_n^{\text{Bohr}} \cdot \phi \cdot \exp\left(-\xi \frac{n^2}{D_f}\right) \tag{ML-Eq.~3.2}
	\end{equation}
	\textbf{Prediction}: 3d state at $n=6$: $\Delta E = -0.00061$ eV ($\sim$1.5$\times$10$^{11}$ Hz), testable via 2-photon spectroscopy (IYQ 2026+).
	
	\subsection{Dirac Equation: Spin-Dependent Corrections}
	
	T0-Original (Section 4.2) modifies Dirac as:
	\begin{equation}
		\left[i\gamma^\mu \left(\partial_\mu + \frac{\xi}{E_{\text{Pl}}} \Gamma_\mu^{(T)}\right) - m\right]\psi = 0 \tag{T0-Orig Eq.~4.2.1}
	\end{equation}
	ML simulations (g-2 anomaly fits) reveal $\xi$-enhancement for heavy leptons.
	
	\textbf{ML-Extended g-Factor}:
	\begin{equation}
		g_{\text{factor}}^{\text{T0,ext}} = 2 + \frac{\alpha}{\pi} + \xi \left(\frac{m}{M_{\text{Pl}}}\right)^2 \cdot \exp\left(-\xi \frac{m}{m_e}\right) \tag{ML-Eq.~3.3}
	\end{equation}
	\textbf{Impact}: Muon $a_\mu$: $\Delta=-0.39\%$ (vs.~Fermilab 2021); electron $a_e$: $\Delta=+0.15\%$ (Schwinger term only).
	
	\section{ML-Derived Neutrino Physics}
	
	\subsection{$\xi^2$-Suppression Mechanism}
	
	T0-Original introduces $\xi^2$ via photon analogy; ML validates via PMNS fits.
	
	\textbf{QFT-Neutrino Propagator}:
	\begin{equation}
		(\Delta m_{ij}^2)^{\text{T0}} \propto \xi^2 \frac{\langle\delta E\rangle}{E_0^2} \approx 10^{-5} \text{ eV}^2 \tag{ML-Eq.~4.1}
	\end{equation}
	\textbf{Hierarchy via $\phi$-Scaling}:
	\begin{align}
		\Delta m_{21}^2 &= \xi^2 \cdot (E_0 / \phi)^2 = 7.52\times10^{-5} \text{ eV}^2 \quad (\Delta=0.4\% \text{ to NuFit}) \tag{ML-Eq.~4.2a} \\
		\Delta m_{31}^2 &= \xi^2 \cdot E_0^2 \cdot \phi = 2.52\times10^{-3} \text{ eV}^2 \quad (\Delta=0.28\%) \tag{ML-Eq.~4.2b}
	\end{align}
	The stated numerical values do not follow from the formulas: ML-Eq.~4.2a/b imply $\Delta m_{31}^2/\Delta m_{21}^2 = \phi^3 = 4.24$, the stated values give $33.5$; with $E_0 = 7.398$~MeV one gets $\Delta m_{21}^2 = 3.7\times10^{5}$~eV$^2$; $7.52\times10^{-5}$~eV$^2$ would require $E_0 \approx 105$~eV, and then $\Delta m_{31}^2 = 3.2\times10^{-4}$~eV$^2$.
	
	\subsection{DUNE Predictions}
	
	\textbf{T0-Oscillation Probability}:
	\begin{equation}
		P(\nu_\mu \to \nu_e)^{\text{T0}} = \sin^2(2\theta_{13}) \sin^2\left(\frac{\Delta m_{31}^2 L}{4E}\right) \cdot \left(1 - \xi \frac{(L/\lambda)^2}{D_f}\right) + \delta E \tag{ML-Eq.~4.3}
	\end{equation}
	\textbf{CP-Violation}: T0 predicts $\delta_{\text{CP}} = 185^\circ \pm 15^\circ$ (NO, $\Delta=13\%$ to NuFit central $212^\circ$)---3$\sigma$ detectable in 3.5 years.
	
	
	
\textbf{Sterile Neutrinos:} The Galois algebra predicts no sterile neutrinos (Doc.~343 Theorem~D$'''$): the Majorana
layer $\{f_1,f_2\}$ is fully occupied by $\nu_1,\nu_2$; $\nu_3$ in orbit $O_4$
is a Dirac neutrino. Authoritative: Doc.~340 and Doc.~343.

\textbf{Testability}: First DUNE runs (2026): prediction $\chi^2$/DOF $<1.1$ for T0-PMNS; $\xi^3$-suppression ($\Delta P<10^{-3}$), not to be interpreted as a sterile signal.
	
	\section{Unified Fractal Framework Across Scales}
	
	\subsection{Universal Damping Pattern}
	
	ML-divergences (QM $n=6$: 44\%, Bell $5\pi/4$: 12.3\%, QFT $\mu=10$ GeV: 0.03\%) converge to:
	
	\begin{tcolorbox}[colback=orange!5!white,colframe=orange!75!black,title={Unified T0 Fractal Law}]
		\begin{equation}
			\mathcal{O}^{\text{T0}}(\text{scale}) = \mathcal{O}^{\text{std}}(\text{scale}) \cdot \exp\left(-\xi \frac{(\text{scale}/\text{scale}_0)^2}{D_f}\right) \tag{ML-Eq.~5.1}
		\end{equation}
		\textbf{Applications}:
		\begin{itemize}
			\item QM: scale $= n$ (Rydberg), scale$_0=1$
			\item Bell: scale $= \Delta\theta/\pi$, scale$_0=1$
			\item QFT: scale $= \ln(\mu/\Lambda_{\text{QCD}})$, scale$_0=1$
		\end{itemize}
	\end{tcolorbox}
	
	\subsection{Emergent Non-Perturbative Structure}
	
	\textbf{Perturbative Expansion} (Taylor of ML-Eq.~5.1):
	\begin{equation}
		\mathcal{O}^{\text{T0}} \approx \mathcal{O}^{\text{std}} \left(1 - \frac{\xi}{D_f} \left(\frac{\text{scale}}{\text{scale}_0}\right)^2 + \mathcal{O}(\xi^2)\right) \tag{ML-Eq.~5.2}
	\end{equation}
	\textbf{Insight}: Linear $\xi$-corrections (T0-Original) are $\mathcal{O}(\xi)$-accurate; ML reveals $\mathcal{O}(\xi \cdot \text{scale}^2)$ at boundaries.
	
	\textbf{Comparison Table}:
	\begin{table}[htbp]
		\centering
		\resizebox{\textwidth}{!}{
			\begin{tabular}{lccc}
				\toprule
				\textbf{Domain} & \textbf{T0-Original $\Delta$} & \textbf{ML-Extended $\Delta$} & \textbf{Improvement} \\
				\midrule
				QM (n=6) & 44\% (divergent) & 0.16\% & +99.6\% \\
				Bell ($5\pi/4$) & 12.3\% & 0.09\% & +99.3\% \\
				QFT ($\mu=10$ GeV) & 0.03\% & 0.008\% & +73\% \\
				\bottomrule
			\end{tabular}
		}
		\caption{ML-Extension Impact Across T0 Applications}
	\end{table}
	
	\subsection{$\phi$-Scaling}
	
	\textbf{Critical Finding}: ML NNs learn prescribed $\phi$-hierarchies exactly (0\% training $\Delta$). They agree with the measured values only in part:
	\begin{itemize}
		\item Masses: a generation ratio $\phi^2 = 2.618$ does not match $m_\mu/m_e = 206.77$; the lepton masses follow the lepton ladder (Section~8.1)
		\item Neutrinos: $\Delta m_{31}^2/\Delta m_{21}^2 = 2.513\times10^{-3}/7.49\times10^{-5} = 33.6$ against $\phi^3 = 4.236$ -- no agreement
		\item Energies: $E_{n,\text{gen}=1} / E_{n,\text{gen}=0} = \phi$ (Rydberg)
	\end{itemize}
	\textbf{Conclusion}: $\phi$-scaling is set geometrically, not ML-emergent; as a mass and neutrino ratio it is not confirmed by the measured values.
	
	\section{Experimental Roadmap}
	
	\subsection{Immediate Tests}
	
	\subsubsection{Loophole-Free Bell Tests}
	
	\textbf{Target}: 100-qubit systems (IBM/Google); T0 predicts:
	\begin{equation}
		\text{CHSH}(N=100) = 2.8278 \pm 0.0001 \quad (\Delta \sim 0.02\%) \tag{ML-Eq.~6.1}
	\end{equation}
	\textbf{Signature}: Deviation from Tsirelson bound ($2.8284$) at $3\sigma$ ($\sim300$ runs).
	
	\subsubsection{Rydberg Spectroscopy}
	
	\textbf{Target}: n=6--20 hydrogen transitions (MPD upgrades); T0 predicts:
	\begin{itemize}
		\item $n=6$: $\Delta E = -6.1\times10^{-4}$ eV ($\sim$1.5$\times$10$^{11}$ Hz)
		\item $n=20$: $\Delta E = -6\times10^{-4}$ eV (cumulative from $n=1$)
	\end{itemize}
	\textbf{Precision}: 2-photon spectroscopy ($\sim$1 kHz resolution); T0 detectable at 5$\sigma$.
	
	\subsection{Medium-Term Tests}
	
	\subsubsection{DUNE First Data}
	
	\textbf{Target}: $\nu_\mu \to \nu_e$ appearance (L=1300 km, E=1--5 GeV); T0 predicts:
	\begin{equation}
		P(\nu_\mu \to \nu_e) = 0.081 \pm 0.002 \quad \text{at } E=3 \text{ GeV} \tag{ML-Eq.~6.2}
	\end{equation}
	\textbf{CP-Violation}: $\delta_{\text{CP}} = 185^\circ$ testable at 3.2$\sigma$ in 3.5 years (vs.~3.0$\sigma$ Standard).
	
	\subsubsection{HL-LHC Higgs Couplings}
	
	\textbf{Target}: $\lambda(\mu=125$ GeV) via $t\bar{t}H$ production; T0 predicts:
	\begin{equation}
		\lambda^{\text{T0}} = 1.0002 \pm 0.0001 \tag{ML-Eq.~6.3}
	\end{equation}
	\textbf{Measurement}: The HL-LHC (3000 fb$^{-1}$) reaches a precision of a few \% for $t\bar{t}H$ and about 30--50\% for the Higgs self-coupling (SM: $\lambda \approx 0.13$; $\lambda^{\text{T0}}$ above is normalized to the SM value); a deviation of $2\times10^{-4}$ cannot be resolved.
	
	\subsection{Long-Term}
	
	\subsubsection{Gravitational Wave T0 Signatures}
	
	\textbf{LIGO-India/ET}: Frequency-dependent corrections:
	\begin{equation}
		h_{\text{T0}}(f) = h_{\text{GR}}(f) \left(1 + \xi \left(\frac{f}{f_{\text{Pl}}}\right)^2\right) \tag{T0-Orig Eq.~8.1.2}
	\end{equation}
	\textbf{Detectability}: Binary mergers at $f\sim100$ Hz: $\Delta h/h \sim 4\times10^{-87}$.
	
	\subsubsection{T0 Quantum Computer Prototype}
	
	\textbf{Target}: Deterministic QC with time-field control; T0 predicts:
	\begin{equation}
		\epsilon_{\text{gate}}^{\text{T0}} = \epsilon_{\text{std}} \cdot \left(1 - \xi \frac{E_{\text{gate}}}{E_{\text{Pl}}}\right) \sim 10^{-5} \tag{T0-Orig Eq.~5.2.1}
	\end{equation}
	\textbf{Benchmark}: Shor's algorithm with $P_{\text{success}}^{\text{T0}} = P_{\text{std}} \cdot (1 + \xi\sqrt{n})$ (n=RSA-2048: +0.6\% boost).
	
	\section{Critical Evaluation and Philosophical Implications}
	
	\subsection{ML's Role: Calibration vs.~Discovery}
	
	\textbf{Key Insight}: ML does \textit{not} replace T0's geometric core---it \textit{reveals} non-perturbative boundaries.
	
	\begin{tcolorbox}[colback=red!5!white,colframe=red!75!black,title={ML Limitations in T0}]
		\textbf{What ML Achieves}:
		\begin{itemize}
			\item Identifies divergences ($\Delta>10\%$) signaling missing terms
			\item Validates $\phi$-scaling (0\% training error)
		\end{itemize}
		\textbf{What ML Cannot Do}:
		\begin{itemize}
			\item Generate $\phi$-hierarchies (purely geometric)
			\item Predict new physics without T0 framework
			\item Replace harmonic formulas (ML gains $<3\%$)
		\end{itemize}
	\end{tcolorbox}
	
	\textbf{Conclusion}: T0 remains parameter-free; ML is a \textit{precision tool}, not a theory builder.
	
	\subsection{Determinism vs.~Practical Unpredictability}
	
	T0-Original (Section 9.1) claims determinism via time fields. \textbf{ML Caveat}:
	\begin{itemize}
		\item \textbf{Sensitivity}: $\xi$-dynamics chaotic at Planck scale ($\Delta E \sim E_{\text{Pl}}$)
		\item \textbf{Computability}: Fractal terms ($\exp(-\xi n^2)$) require infinite precision for $n\to\infty$
		\item \textbf{Effective Randomness}: Bell outcomes deterministic in principle, but computationally inaccessible
	\end{itemize}
	\textbf{Philosophical Stance}: T0 restores ontological determinism, but preserves epistemic uncertainty---reconciling Einstein's ``God does not play dice'' with Born's probabilistic observations.
	
	
	\subsection{Locality and Bell's Theorem}
	
	T0-Original (Section 6.2) claims local hidden variables via time fields. \textbf{ML Insight}:
	\begin{equation}
		\lambda_{\text{T0}} = \{T_{\text{field},A}(t), T_{\text{field},B}(t), \text{common history}\} \tag{ML-Eq.~7.1}
	\end{equation}
	\textbf{Objection}: Does the theoretical $\text{CHSH}^{\text{T0}}\approx2.8279$ violate Bell's bound (2)?
	
	\textbf{Answer}: Yes. T0 modifies \textit{expectation values}:
	\begin{itemize}
		\item Standard Bell assumes $E(a,b) = \int P(A,B|a,b,\lambda) \cdot A \cdot B \, d\lambda$
		\item T0 adds: $E^{\text{T0}}(a,b) = \int P(\cdots) \cdot A \cdot B \cdot \exp(-\xi f(\lambda)) \, d\lambda$
		\item Result: with a weight $\exp(-\xi f) \leq 1$ a local model stays at $|S| \leq 2$; with $\xi\Delta \ll 1$ the modified bound is also only $|S| \leq 2 + \xi\Delta \approx 2.0001$
	\end{itemize}
	\textbf{Critical Point}: The document's own value $\text{CHSH}^{\text{T0}} = 2.8279$ would require $\xi\Delta = 0.83$, i.e.\ $\Delta \approx 6200$, not a small $\xi$ correction. Loophole-free Bell tests measure $S$ clearly above $2$ (e.g.\ Storz et al.\ 2023: $2.0747 \pm 0.0033$; photonic experiments up to about $2.7$); a local model with this bound is therefore excluded.
	
	\section{Synthesis: The T0-ML Unified Picture}
	
	\subsection{Three-Tier Hierarchy of T0 Theory}
	
	\begin{tcolorbox}[colback=blue!5!white,colframe=blue!75!black,title={T0 Theoretical Structure}]
		\textbf{Tier 1: Geometric Foundation} (Parameter-Free)
		\begin{itemize}
			\item $\xi = 4/30000$ (fractal dimension $D_f=3-\xi$)
			\item $\phi = (1+\sqrt{5})/2$ (golden ratio scaling)
			\item $T_{\text{field}} \cdot E_{\text{field}} = 1$ (time-energy duality)
		\end{itemize}
		
		\textbf{Tier 2: Harmonic Predictions} (1--3\% Precision)
		\begin{itemize}
			\item Masses: $m_i = r_i \cdot \xi^{p_i} \cdot v$ (lepton ladder, Doc.~006; $r = 4/3, 16/5, 25/9$, $p = 3/2, 1, 2/3$, $v = 246$~GeV)
			\item Neutrinos: $\Delta m^2 \propto \xi^2 \cdot \phi^{\text{hierarchy}}$
			\item QM: $E_n = E_n^{\text{Bohr}} \cdot (1 + \xi E_n/E_{\text{Pl}})$
		\end{itemize}
		
		\textbf{Tier 3: ML-Derived Extensions} (0.1--1\% Precision)
		\begin{itemize}
			\item Fractal damping: $\exp(-\xi \cdot \text{scale}^2/D_f)$
			\item QFT loops: Natural cutoff $\Lambda_{\text{T0}} = E_{\text{Pl}}/\xi$
		\end{itemize}
	\end{tcolorbox}
	
	\subsection{Predictive Power Comparison}
	
	\begin{table}[htbp]
		\centering
		\resizebox{\textwidth}{!}{
			\begin{tabular}{lcc}
				\toprule
				\textbf{Observable} & \textbf{SM (Free Params)} & \textbf{T0 Geometric} \\
				\midrule
				Lepton Masses & 3 (fitted) & $\Delta=0.4$--$1.2\%$ \\
				Neutrino $\Delta m^2$ & 2 (fitted) & not derivable from ML-Eq.~4.2 \\
				CHSH (Bell) & N/A (QM: 2.828) & $\Delta=0.02\%$ \\
				Higgs Mass & 1 (fitted) & $\Delta=0.1\%$ \\
				Hydrogen $E_6$ & 0 (QED exact) & $\Delta=0.12\%$ \\
				\midrule
				Total Free Params & $\sim$19 (SM) & 0 ($\xi, \phi$ geometric) \\
				\bottomrule
			\end{tabular}
		}
		\caption{T0 vs.~Standard Model: Predictive Precision}
	\end{table}
	The ladder $m_i = r_i\xi^{p_i}v$ gives $m_e = 0.5050$~MeV ($-1.18\%$), $m_\mu = 104.96$~MeV ($-0.66\%$), $m_\tau = 1783.4$~MeV ($+0.37\%$).
	
	
	\subsection{Appendix: ML Training Details (For Reproducibility)}
	
	\textbf{Bell Correlation NN}:
	\begin{itemize}
		\item Architecture: Input(3: $a, b, \xi$) $\to$ Dense(32, ReLU) $\to$ Dense(16, ReLU) $\to$ Output(1: $E(a,b)$)
		\item Loss: MSE to QM $E=-\cos(a-b)$
		\item Training: 1000 samples ($\Delta\theta \in [0,\pi/2]$), 200 epochs, Adam($\eta=10^{-3}$)
		\item Test: $\Delta\theta \in [\pi/2, 2\pi]$; Divergence at $5\pi/4$: 12.3\%
	\end{itemize}
	
	\textbf{Rydberg Energy NN}:
	\begin{itemize}
		\item Architecture: Input(1: $n$) $\to$ Dense(64, Tanh) $\to$ Dense(32, Tanh) $\to$ Output(1: $E_n$)
		\item Loss: MSE to Bohr $E_n = -13.6/n^2$
		\item Training: $n=1$--5 (5 samples), 500 epochs; Test: $n=6$ diverges (44\%)
		\item Fix: Integrate $\exp(-\xi n^2/D_f)$; Retraining: $\Delta<0.2\%$ for $n=1$--20
	\end{itemize}
	
	
	\section{Comparative Table: T0-Original vs.~T0-ML}
	
	\section*{Comparison Table}
	\begin{longtable}{p{2.4cm}p{4.0cm}p{4.0cm}}
		\toprule
		\textbf{Aspect} & \textbf{T0-Original (2025)} & \textbf{T0-ML Addendum (2025)} \\
		\midrule
		\endfirsthead
		\toprule
		\textbf{Aspect} & \textbf{T0-Original} & \textbf{T0-ML Addendum} \\
		\midrule
		\endhead
		
		Bell CHSH & $2 + \xi\Delta_{\text{T0}}$ (qualitative) & $2.8279$ (N=73, theoret.; HW: $2.74$) \\
		QM Hydrogen & $E_n(1+\xi E_n/E_{\text{Pl}})$ & $E_n \cdot \phi^{\text{gen}} \cdot \exp(-\xi n^2/D_f)$ \\
		Neutrino Mass & $\xi^2$-suppression (concept) & $\Delta m_{21}^2=7.52\times10^{-5}$ eV$^2$ (does not follow from ML-Eq.~4.2a) \\
		ML Role & Not discussed & Precision tool (0.1--3\% gain) \\
		Testability & Qualitative predictions & Quantitative (DUNE $\delta_{\text{CP}}=185^\circ$) \\
		Fractal Terms & Implied in $D_f$ & Explicit $\exp(-\xi \cdot \text{scale}^2/D_f)$ \\
		Precision & $\sim$1--3\% (harmonic) & $\sim$0.1--1\% (ML-extended) \\
		\bottomrule
		\caption{Comprehensive Comparison: T0-Original vs.~ML Extensions}
	\end{longtable}
	
	\section{Glossary of Key Terms}
	
	\begin{description}
		\item[Fractal Damping] $\exp(-\xi \cdot \text{scale}^2/D_f)$ correction stabilizing divergences at boundary scales (high $n$, angles, $\mu$).
		\item[$\phi$-Scaling] Golden ratio hierarchies ($\phi^{\text{gen}}$), e.g.\ in Rydberg energies---learned exactly by ML (0\% error); not confirmed as a mass ratio.
		\item[ML Divergence] NN prediction error $>10\%$ at test boundaries, signaling missing physics (emergent terms).
		\item[T0-Original] Base document (020\_T0\_QM-QFT-RT\_En.pdf) establishing time-energy duality and QFT framework.
		\item[Loophole-Free] Bell tests with $>$95\% detection efficiency, excluding local hidden variable explanations.
	\end{description}
	
	\begin{thebibliography}{99}
		
		\bibitem{pascher_t0_qft_2025}
		Pascher, J. (2025). \textit{T0 Quantum Field Theory: Complete Extension --- QFT, QM and Quantum Computers}.
		T0-Original Document (020\_T0\_QM-QFT-RT\_En.pdf).
		
		\bibitem{pascher_bell_ml_2025}
		Pascher, J. (2025). \textit{T0-Theorie: Erweiterung auf Bell-Tests --- ML-Simulationen}.
		023\_Bell\_En.pdf, November 2025.
		
		\bibitem{pascher_qm_summary_2025}
		Pascher, J. (2025). \textit{T0-Theorie: Zusammenfassung der Erkenntnisse}.
		035\_QM\_En.pdf, Stand November 03, 2025.
		
		
		\bibitem{mpd_hydrogen_2025}
		MPD Collaboration (2025). \textit{Metrology for Precise Determination of Hydrogen Energy Levels}.
		arXiv:2403.14021v2 [physics.atom-ph], May 2025.
		
		\bibitem{nufit_2024}
		Esteban, I., et al. (2024). \textit{NuFit 6.0: Updated Global Analysis of Neutrino Oscillations}.
		\url{http://www.nu-fit.org}, September 2024.
		
		\bibitem{dune_2025}
		DUNE Collaboration (2025). \textit{Deep Underground Neutrino Experiment: Physics Prospects}.
		NuFact 2025 Conference Proceedings.
		
		\bibitem{particle_data_group_2024}
		Particle Data Group (2024). \textit{Review of Particle Physics}.
		Prog. Theor. Exp. Phys. \textbf{2024}, 083C01.
		
		\bibitem{iyq_2025}
		International Year of Quantum (2025). \textit{About IYQ}.
		\url{https://quantum2025.org/about/}
		
	\end{thebibliography}
	